\chapter{Wie funktioniert ein Kampf?}

Ein Kampf ist eine besondere Szene, in der die Zeit in Runden und Züge aufgeteilt wird.

\section{Die Kampfrunde}

In jeder Runde handeln zuerst alle Helden, danach alle Gegner. Die Helden bestimmen ihre Reihenfolge selbst. Sobald alle Gegner vollständig aktiviert wurden, endet die Runde und die nächste beginnt.

\subsection{Der Zug eines Helden}

Ein Held besitzt in seinem Zug normalerweise:

\begin{itemize}
  \item eine \textbf{Standardaktion (\standardAction)},
  \item eine \textbf{Bewegungsaktion (\moveAction)},
  \item beliebig viele \textbf{freie Aktionen (\freeAction)}.
\end{itemize}

Mit einer \standardAction{} kann ein Held beispielsweise angreifen, einen Zauber wirken oder eine besondere Fähigkeit einsetzen.

Mit einer \moveAction{} kann er sich bewegen oder seine Ausrüstung neu anordnen. Seine \standardAction{} darf er stattdessen als zweite \moveAction{} verwenden.

\freeAction{} sind kurze Handlungen wie etwas rufen, auf etwas zeigen oder einen Gegenstand fallen lassen. Sie dürfen keine umfangreiche Handlung ersetzen.

\section{Bewegung, Nähe und Entfernung}

Figuren können zueinander \textbf{benachbart} oder \textbf{entfernt} sein. Eine Figur kann gleichzeitig zu mehreren Figuren benachbart sein.

Wie viele Bewegungen nötig oder mit einer \moveAction{} möglich sind, ergibt sich aus der Battlemap oder aus der Beschreibung der Spielleitung. Hindernisse und die konkrete Situation können eine Bewegung erschweren oder mehrere Bewegungen erforderlich machen.

Ein Held darf sich von einem benachbarten Gegner entfernen, ohne dass dies automatisch eine weitere Folge auslöst. Gegnerfähigkeiten können davon abweichen.

\section{Ausrüstung im Kampf wechseln}

Mit einer \moveAction{} darf ein Held seine ausgerüsteten und verstauten Gegenstände beliebig neu anordnen, solange danach alle Ausrüstungsplätze korrekt belegt sind.

Der Gegenstand am Körper-Ausrüstungsplatz ist davon ausgenommen. Eine Körperrüstung anzulegen, abzulegen oder zu wechseln, benötigt eine eigene \standardAction.

\section{Einen Gegner angreifen}

Ein Angriff benötigt eine \standardAction. Die gewählte Angriffsoption bestimmt die Angriffsfertigkeit, die Reichweite, die Schadensformel und mögliche besondere Effekte.

\begin{itemize}
  \item \textbf{Schlagen} kann benachbarte Gegner als Ziel haben.
  \item \textbf{Schießen} kann entfernte Gegner als Ziel haben. Ein Held darf auch schießen, während ein Gegner zu ihm benachbart ist; das Ziel des Angriffs muss dennoch entfernt sein.
  \item \textbf{Zaubern} kann benachbarte oder entfernte Gegner als Ziel haben, sofern der Zauber nichts anderes bestimmt.
\end{itemize}

\subsection{Treffen}

Der Spieler legt mit der angegebenen Angriffsfertigkeit eine Angriffsprobe ab. Zeigt sie keinen \star, verfehlt der Angriff. Jeder erzielte \star{} wird anschließend nach der Schadensformel der Angriffsoption in Schadenswürfel umgerechnet.

Helden erhalten Vorteil auf Angriffsproben gegen am Boden liegende Gegner.

\subsection{Schaden machen}

Es gibt drei Arten von Schadenswürfeln:

\begin{itemize}
  \item den \textbf{orangen Schadenswürfel (\orangeDmgDie)} für leichten Schaden,
  \item den \textbf{roten Schadenswürfel (\redDmgDie)} für mittleren Schaden,
  \item den \textbf{schwarzen Schadenswürfel (\blackDmgDie)} für schweren Schaden.
\end{itemize}

Die Schadensformel einer Angriffsoption bestimmt, wie viele Schadenswürfel jeder erzielte \star{} erzeugt.

\textbf{Beispiel:} Eine Angriffsoption erzeugt pro \star{} zwei \orangeDmgDie. Erzielt Kiera 3\star, würfelt sie 6\orangeDmgDie.

Die Schadenswürfel werden geworfen und ihr Schaden wird addiert. Danach wird der Schaden durch Schild-Symbole verringert.

Das \textbf{Schild-Symbol (\armor)} beschreibt Schutz durch Rüstung, Schilde oder andere Effekte. Jedes \armor{} verhindert 1 Schaden. \textbf{Rüstungsdurchdringung (\armorPierce)} ignoriert die angegebene Anzahl an \armor{} des Ziels.

Der verbleibende Schaden lässt die Figur entsprechend viele \textbf{Lebenspunkte verlieren (\lifeLoss)}.

\section{Wenn Gegner angreifen}

Gegner legen keine Angriffsprobe ab. Ihr Profil gibt unmittelbar die verwendeten Schadenswürfel an.

Bevor diese Würfel geworfen werden, darf der angegriffene Held gegebenenfalls \textbf{ausweichen (\evade)}. Welche Würfel er für seine Ausweichprobe verwendet, ergibt sich aus seiner Ausrüstung.

Jeder \star{} der Ausweichprobe entfernt einen beliebigen Schadenswürfel des Gegners. Entfernte Schadenswürfel werden nicht geworfen.

Ein am Boden liegender Held hat Nachteil auf seine Ausweichprobe. Kann er aufgrund seiner Ausrüstung nicht ausweichen, entfällt die Probe vollständig.

Die verbleibenden Schadenswürfel werden geworfen. Anschließend verhindert jedes \armor{} des Helden 1 Schaden. Für den übrigen Schaden verliert der Held entsprechend viele Lebenspunkte \lifeLoss.

\section{Kampfunfähig oder besiegt}

Lebenspunkte können niemals unter 0 sinken.

Sinkt ein Gegner auf 0 Lebenspunkte, ist er besiegt.

Sinkt ein Held auf 0 Lebenspunkte, ist er kampfunfähig und kann nicht handeln. Wird er geheilt und erhält mindestens 1 \textbf{Lebenspunkt (\lifeGain)}, ist er wieder kampfbereit, erhält aber den Status \emph{Am Boden}. Er muss daher zunächst entsprechend den Statusregeln aufstehen.

Wenn alle Helden kampfunfähig sind, sind die Helden besiegt. Das bedeutet nicht automatisch, dass sie sterben.

\section{Wann endet ein Kampf?}

Ein Kampf endet, wenn die gefährliche Auseinandersetzung vorbei ist, beispielsweise weil alle Gegner besiegt sind, eine Seite flieht oder aufgibt oder das Ziel des Kampfes erreicht wurde.

Wenn der Kampf endet, wird wieder normal weitererzählt.
